% ============================================================
%  BUDGET PRÉVISIONNEL — Phase 2 : Planification
%  Time'Eats — Cellule PMO
% ============================================================
\documentclass[11pt,a4paper,oneside]{report}
\usepackage{../style/consulting}
\hypersetup{pdftitle={Budget Prévisionnel — Time'Eats PMO}}
\pagestyle{fancy}\fancyhf{}
\renewcommand{\headrulewidth}{0pt}
\fancyhead[L]{\begin{tikzpicture}[remember picture,overlay]
  \draw[cnNavy,line width=0.5pt](0,0)--(\textwidth,0);
\end{tikzpicture}\vspace{2pt}\timeeatslogo\hspace{4pt}\small\textcolor{cnSlate}{Time'Eats \textbf{·} Budget Prévisionnel — \textit{[Nom projet]}}}
\fancyhead[R]{\small\textcolor{cnSlate}{\thepage}}
\fancyfoot[C]{\tiny\textcolor{cnSlate}{Document confidentiel — Cellule PMO — CESI MAALSI 2026}}
\begin{document}\small

\begin{center}
{\LARGE\bfseries\color{cnNavy} Budget Prévisionnel}\\[4pt]
\textcolor{cnOrange}{\rule{10cm}{1pt}}\\[4pt]
{\large\color{cnSlate} Phase 2 — Planification \quad|\quad Projet : \textit{[Nom du projet]}}
\end{center}

\vspace{6pt}
\renewcommand{\arraystretch}{1.4}
\begin{tabularx}{\textwidth}{L{3.5cm} Y L{3.5cm} Y}
\toprule
\rowcolor{cnNavy}\th{Champ} & \th{Valeur} & \th{Champ} & \th{Valeur} \\
\midrule
Budget total autorisé & \textit{[XXX K€]} & Version   & 1.0 \\
\rowcolor{cnIce}
Rédacteur             & \textit{[Prénom NOM]} & Date  & \textit{[JJ/MM/AAAA]} \\
\bottomrule
\end{tabularx}

\section*{1 — Décomposition budgétaire par poste}

\renewcommand{\arraystretch}{1.4}
\begin{tabularx}{\textwidth}{L{5cm} C{2.5cm} C{2.5cm} C{2.5cm} Y}
\toprule
\rowcolor{cnNavy}\th{Poste de dépense} & \th{Type} & \th{Montant (K€)} & \th{\% budget} & \th{Commentaire} \\
\midrule
Licences / SaaS        & OPEX & \textit{[X]} & \textit{[X\%]} & \textit{[Détail]} \\
\rowcolor{cnIce}
Intégration ESN        & OPEX\textsuperscript{(*)} & \textit{[X]} & \textit{[X\%]} & Forfait --- prestations de services \\
Ressources internes    & Interne & \textit{[X]} & \textit{[X\%]} & TJM estimé \\
\rowcolor{cnIce}
Infrastructure / Hébergement & CAPEX & \textit{[X]} & \textit{[X\%]} & \textit{[Détail]} \\
Formation utilisateurs & OPEX & \textit{[X]} & \textit{[X\%]} & \textit{[Détail]} \\
\rowcolor{cnIce}
Tests et recette       & Interne & \textit{[X]} & \textit{[X\%]} & \textit{[Détail]} \\
Conduite du changement & OPEX & \textit{[X]} & \textit{[X\%]} & \textit{[Détail]} \\
\rowcolor{cnIce}
Réserve pour aléas     & Réserve & \textit{[X]} & 10\% & Non affecté au départ \\
\midrule
\rowcolor{cnNavy!10}\textbf{TOTAL} & & \textbf{\textit{[X K€]}} & \textbf{100\%} & \\
\bottomrule
\end{tabularx}

\vspace{4pt}
{\footnotesize
\textsuperscript{(*)} Les j/h ESN sont par défaut comptabilisés en \textbf{OPEX} (prestations de services consommées dans l'exercice). Une fraction peut être immobilisée en \textbf{CAPEX} si elle correspond à du développement spécifique répondant aux critères IAS 38 (faisabilité technique, intention d'achever, capacité à utiliser, génération d'avantages économiques futurs, ressources adéquates, mesure fiable des coûts). \textit{Arbitrage à formaliser avec la DAF projet par projet et à tracer dans le présent budget.}
}

\section*{2 — Répartition temporelle (échelonnement)}

\renewcommand{\arraystretch}{1.3}
\begin{tabularx}{\textwidth}{L{5cm} C{2cm} C{2cm} C{2cm} C{2cm} C{2cm}}
\toprule
\rowcolor{cnNavy}\th{Poste} & \th{M1--M2} & \th{M3--M4} & \th{M5--M6} & \th{M7--M8} & \th{M9+} \\
\midrule
Licences / SaaS      & \textit{[X]} & \textit{[X]} & \textit{[X]} & \textit{[X]} & \textit{[X]} \\
\rowcolor{cnIce}
Intégration ESN      & \textit{[X]} & \textit{[X]} & \textit{[X]} & \textit{[X]} & \textit{[X]} \\
Ressources internes  & \textit{[X]} & \textit{[X]} & \textit{[X]} & \textit{[X]} & \textit{[X]} \\
\rowcolor{cnIce}
Formation            & --           & --           & --           & \textit{[X]} & \textit{[X]} \\
\midrule
\rowcolor{cnNavy!10}\textbf{Total mensuel} & \textit{[X]} & \textit{[X]} & \textit{[X]} & \textit{[X]} & \textit{[X]} \\
\bottomrule
\end{tabularx}

\section*{3 — Suivi consommé vs prévu (EVM)}

\renewcommand{\arraystretch}{1.3}
\begin{tabularx}{\textwidth}{L{4cm} C{2.5cm} C{2.5cm} C{2.5cm} Y}
\toprule
\rowcolor{cnNavy}\th{Période} & \th{BAC prévu (K€)} & \th{CBTE (K€)} & \th{CRTE (K€)} & \th{IPC / IPD} \\
\midrule
Mois 1 & \textit{[X]} & \textit{[X]} & \textit{[X]} & \textit{[X.XX / X.XX]} \\
\rowcolor{cnIce}
Mois 2 & \textit{[X]} & \textit{[X]} & \textit{[X]} & \textit{[X.XX / X.XX]} \\
Mois 3 & \textit{[X]} & \textit{[X]} & \textit{[X]} & \textit{[X.XX / X.XX]} \\
\bottomrule
\end{tabularx}

\vspace{4pt}
{\footnotesize
\textbf{BAC} = Budget à Complétion \quad
\textbf{CBTE} = Coût Budgété du Travail Effectué (valeur acquise) \quad
\textbf{CRTE} = Coût Réel du Travail Effectué \quad
\textbf{IPC} = CBTE/CRTE \quad
\textbf{IPD} = CBTE/CBTP
}

\begin{reco}
IPC \textless{} 0.90 ou IPD \textless{} 0.90 → déclencher un \textbf{rapport d'écarts} immédiat et convoquer un COPIL exceptionnel.
\end{reco}

\end{document}
